% ============================================================================
% Fichier : annexes/Ann_B_outils.tex
% Titre   : Référentiel des Outils
% Version : 1.0 -- novembre 2026
% ============================================================================
% Outils mobilisés dans les chapitres et les labs ; les noms commerciaux sont
% des exemples, à vérifier au moment du choix.
% ============================================================================

\chapter{Référentiel des Outils}
\label{ann:outils}

Cette annexe rassemble les outils de la gouvernance et de la gestion des
services mobilisés dans le cours et dans les labs de NDAP SUKULU. Les noms
de logiciels sont des exemples~: les offres évoluent, et un outil ne
remplace ni une décision ni une procédure.

\section*{Gouvernance et conception}

\begin{tableaucours}{L{3.6cm}L{5.4cm}L{3.2cm}}{%
    \tableauHeader{Outil} &
    \tableauHeader{Usage} &
    \tableauHeader{Chapitre, lab}
}
Cascade des objectifs (tableur) & Relier besoins des parties prenantes,
objectifs d'entreprise, d'alignement et objectifs COBIT &
Ch.~\ref{chap:cobit-objectifs}, GCI-LAB-03 \\
Canevas de conception & Appliquer les onze facteurs de conception et fixer
les capacités cibles & Ch.~\ref{chap:cobit-conception}, GCI-LAB-02 \\
Grille d'évaluation de capacité & Coter les activités d'un processus,
niveau par niveau & Ch.~\ref{chap:cobit-capacite}, GCI-LAB-04 \\
Tableau de portefeuille & VAN, coût total de possession, risque, alignement
& Ch.~\ref{chap:alignement}, GCI-LAB-01 \\
Matrice des droits de décision & Qui est consulté, qui décide, pour chaque
type de décision & Ch.~\ref{chap:alignement} \\
\end{tableaucours}

\section*{Gestion des services}

\begin{tableaucours}{L{3.6cm}L{5.4cm}L{3.2cm}}{%
    \tableauHeader{Outil} &
    \tableauHeader{Usage} &
    \tableauHeader{Chapitre, lab}
}
Outil de gestion des services & Tickets d'incidents et de demandes,
priorités, escalade, problèmes & Ch.~\ref{chap:itil-incidents}, GCI-LAB-06 \\
Registre des changements & Demande, évaluation du risque, autorisation,
calendrier, retour arrière & Ch.~\ref{chap:itil-changements}, GCI-LAB-07 \\
CMDB & Éléments de configuration et leurs relations, pour l'analyse
d'impact & Ch.~\ref{chap:itil-changements} \\
Supervision & Mesure de la disponibilité, alertes &
Ch.~\ref{chap:itil-niveaux}, GCI-LAB-08 \\
Matrice RACI & Un seul A par activité~; rôles et non personnes &
Ch.~\ref{chap:processus}, GCI-LAB-09 \\
\end{tableaucours}

\section*{Architecture, livraison et données}

\begin{tableaucours}{L{3.6cm}L{5.4cm}L{3.2cm}}{%
    \tableauHeader{Outil} &
    \tableauHeader{Usage} &
    \tableauHeader{Chapitre}
}
Cadre d'architecture d'entreprise (par exemple TOGAF) & Décrire processus,
données, applications et technologies, et leur évolution &
Ch.~\ref{chap:referentiels} \\
Chaîne de déploiement & Construire, tester et déployer automatiquement,
avec traçabilité & Ch.~\ref{chap:agile-cloud} \\
Catalogue et dictionnaire de données & Propriétaires, définitions,
qualité, conservation & Ch.~\ref{chap:donnees-ia} \\
\end{tableaucours}

\section*{Contrôle et audit}

\begin{tableaucours}{L{3.6cm}L{5.4cm}L{3.2cm}}{%
    \tableauHeader{Outil} &
    \tableauHeader{Usage} &
    \tableauHeader{Chapitre, lab}
}
Extractions et empreintes & Tester les contrôles généraux sur des données
tracées & Ch.~\ref{chap:audit}, GCI-LAB-05 \\
Tableur ou SQL & Rapprochements (départs et comptes actifs, déploiements et
changements autorisés) & Ch.~\ref{chap:audit},
\ref{chap:itil-changements} \\
\end{tableaucours}

\section*{Le bureau de NDAP SUKULU}

Les labs s'exécutent dans un bureau d'analyste virtuel. Les pièces sont
remises dans \texttt{\textasciitilde/preuves}, en lecture seule, avec leur
manifeste d'empreintes~; une copie de travail se trouve dans
\texttt{\textasciitilde/travail}. Le bureau propose un terminal et une suite
bureautique~; l'activité dans le terminal est journalisée.
